\documentclass[11pt]{article}
\usepackage[margin=0.8in]{geometry}
\usepackage{enumitem}
\usepackage{setspace}
\setlist[itemize]{leftmargin=1.15em, itemsep=0.1em, topsep=0.18em}
\setlist[enumerate]{leftmargin=1.35em, itemsep=0.1em, topsep=0.18em}
\setlength{\parindent}{0pt}
\setlength{\parskip}{0.26em}
\setstretch{1.02}
\pagestyle{plain}

\begin{document}

\begin{center}
{\Large\bfseries Industrial Economics}\\[0.3em]
Sample syllabus\\[0.15em]
Ke Cheng HSU\\
Department of Economics, National Chengchi University\\
Undergraduate elective, 3 credits, one semester, taught in English
\end{center}

\textbf{Place in the curriculum.}
This is a one-semester undergraduate elective. It looks at an industry from three sides: the market, the firms, and the regulator. The market weeks cover structure, market power, and where that power comes from. The firm weeks cover pricing, quality, and reputation. The policy weeks cover welfare, collusion, and vertical contracts. A later elective can take up horizontal mergers, research and development, and market intermediation, which are Chapters 15, 18, 19, 22, and 23 of the textbook and are not in this semester.

\textbf{Prerequisites.}
Intermediate microeconomics and basic calculus. Game theory is taught in the course, in the weeks that use it.

\textbf{Textbook.}
Belleflamme, Paul, and Martin Peitz. \emph{Industrial Organization: Markets and Strategies}. 2nd edition. Cambridge University Press, 2015. An earlier edition is acceptable for the models. Chapter numbers below are for the second edition. Appendix A is the game-theory appendix. Problem sets use the exercises at the end of the assigned chapters.

\textbf{How the course runs.}
Three hours a week. Two hours are lecture. The third hour is a problem session for the harder questions on the current set. Six problem sets, posted every other week and due the following week. Students may discuss the setup. Each student hands in their own answers. A serious attempt at the setup earns credit when the last algebra step is wrong. Solutions are posted after the due date.

\textbf{Learning outcomes.}
By the end of the semester, students should be able to
\begin{itemize}
\item measure market power and relate it to the elasticity of demand,
\item solve a static oligopoly model and a simple dynamic one, and say how a cost change moves the equilibrium,
\item explain product differentiation, consumer inertia, and price discrimination as sources of market power,
\item identify a welfare loss from collusion or from double marginalization, and describe one contractual remedy.
\end{itemize}

\textbf{Assessment.}
Problem sets 20 percent. The lowest set is dropped. Midterm examination 30 percent, in class, covering weeks 1--6. A documented absence on the midterm moves that weight to the final. Final examination 50 percent, cumulative. Examinations use the same kind of question as the problem sets. Closed book. A calculator is allowed.

\textbf{Weekly plan.}
\begin{enumerate}
\item Market power and welfare. A review of perfect competition and monopoly pricing. Chapter 2.
\item Static games, as needed for oligopoly. Appendix A, selected sections.
\item Static oligopoly. Cournot and Bertrand. Chapter 3.
\item Dynamic games. Appendix A, selected sections, and the opening of Chapter 4.
\item Dynamic oligopoly and business strategies. Chapter 4.
\item Entry, accommodation, and deterrence, including capacity. Selected sections of Chapter 16.
\item Midterm examination on weeks 1--6.
\item Repeated games and the logic of tacit collusion. Chapter 14.
\item What makes collusion easier or harder, and how a leniency program changes the incentive to report a cartel. The rest of Chapter 14.
\item Product differentiation. Chapter 5.
\item Consumer inertia and switching costs. Chapter 7.
\item Group pricing and personalized pricing. Chapter 8.
\item Menu pricing. Chapter 9. Experience goods and reputation, from Chapter 13, if the week allows.
\item Vertically related markets. Double marginalization and a two-part tariff that restores the integrated quantity. Exclusive dealing. Chapter 17.
\item Markets with network goods. Network effects and compatibility. Chapters 20 and 21.
\item One published paper on pricing, collusion, or a vertical contract. The class reads the question, the data, and the result.
\item Review.
\item Final examination.
\end{enumerate}

\end{document}
